\documentclass[main.tex]{subfiles}


\begin{document}

%from pagenumber to up
% \phantomsection 
% \hypertarget{toc}{}

\pagestyle{empty}

\begin{NoHyper} % отключить клик-ссылки

% номер секции вручную
\setcounter{section}{18
}
\addtocounter{section}{-1} 

\section{Сопротивление}


\textbf{Сопротивление} ($R$ [Ом])~--- это характеристика проводника, показывающая его способность препятствовать прохождению тока через него. Например, если проводник~А подключить к батарейке, а затем (отсоединив проводник~А от батарейки) проводник~Б подключить к этой же батарейке, то сопротивление будет \emph{больше} у того проводника, в котором \emph{сила тока} была \emph{меньше}.   

Пусть имеются три медных провода (рис.\,\ref{fig:p118}).

    \begin{figure}[H]
    \centering
\begin{asy}
settings.outformat="png";
// settings.prc=true;
//settings.render=15;

import three;
import texcolors;
import x11colors;
import solids;
texpreamble("\usepackage[T2A]{fontenc}\usepackage[utf8]{inputenc}
\usepackage{mathtext}\usepackage[russian]{babel}");
size(6cm,0);
currentprojection =
perspective((10,15,10));


//styles
///////////////////////////////////////////
DefaultHead3.size=new real(pen p=currentpen) {return 3mm;};
defaultpen(fontsize(11pt));
defaultpen(linewidth(0.4pt));
pen thick=black+2*linewidth(currentpen);
/////////////////////////////////////////


//axes
//draw(O--2.5X,L=Label("$x$",EndPoint,WSW));
//draw(O--2.3Y,L=Label("$y$",EndPoint,ESE));
transform3 t=shift(-2,-2,0);


real a=4;


/////////////////////////////////////////////
//cylinder1
triple pO=(0,-4,0);
revolution cyl=cylinder(pO,.1,8,Y);
draw(surface(cyl),orange);
draw(shift(0,a,0)*rotate(-90,X)*scale3(.1)*unitdisk,orange);
draw(shift(0,-a,0)*rotate(-90,X)*scale3(.1)*unitdisk,orange);
path3 l=(.1,-a,0)--(.1,a,0);
label("А",(0,-4,.1),N);


//cylinder2
triple pO=(0,-4,0);
revolution cyl=cylinder(pO,.5,8,Y);
draw(shift(-1.6,0,0)*surface(cyl),orange);
draw(shift(-1.6,a,0)*rotate(-90,X)*scale3(.5)*unitdisk,orange);
draw(shift(-1.6,-a,0)*rotate(-90,X)*scale3(.5)*unitdisk,orange);
path3 l=(.1,-a,0)--(.1,a,0);
label("Б",(-1.6,-4,.5),N);


//cylinder3
triple pO=(0,-4,0);
revolution cyl=cylinder(pO,.5,4,Y);
draw(shift(-3.6,2,0)*surface(cyl),orange);
draw(shift(-3.6,-2,0)*rotate(-90,X)*scale3(.5)*unitdisk,orange);
draw(shift(-3.6,2,0)*rotate(-90,X)*scale3(.5)*unitdisk,orange);
path3 l=(.1,-a,0)--(.1,a,0);
label("В",(-3.6,-2,.5),N);


\end{asy}
\caption{Три проводника}
\label{fig:p118}
\end{figure}

Провода~А, Б и~В имеют сопротивления~$R_\text{А}$, $R_\text{Б}$ и~$R_\text{В}$. Опыт показывает, что их сопротивления связаны так: $R_\text{А}>R_\text{Б}>R_\text{В}$. Сравнивать сопротивления тел из одинакового материала можно, используя следующую аналогию: чем <<проще>> протечь воображаемой жидкости через тело, представляемое как <<труба>> (то есть чем шире и короче <<труба>>), тем меньше сопротивление тела.  

\textbf{Сопротивление проводника} вычисляется по формуле:
\begin{equation}
    R=\frac{\rho_\text{э}l}{S},
\end{equation}
где $\rho_\text{э}$~--- \emph{удельное сопротивление} вещества (см.~справочные таблицы), $l$~--- длина проводника, $S$~--- площадь поперечного сечения проводника\footnote{Или площадь поперечной грани (торца) тела, из которой предполагается <<вытекает>> ток.}. 


Сопротивление прохождению тока, например, в металлах объясняют тем, что свободные электроны, совершая направленное движение, сталкиваются с ионами металла, колеблющимися вблизи своих положений равновесия (рис.\,\ref{fig:p120}).

    \begin{figure}[H]
    \centering

    \begin{tikzpicture}[font=\small,            line cap=round,                             >={Stealth[inset=0pt,round,length=3mm, angle'=20]},vec/.style={>={Stealth[inset=0pt,round,length=3.5mm, angle'=20]}}]


% \tikzset{atom/.pic={
%   \begin{scope}[shift={({rnd*360}:.05)}]
%   \shade[ball color=red] (0,0) circle (\dn);
%   % \shade[ball color=NavyBlue!70,rotate={rnd*360}] ([shift={(180:{.4 cm})}]0,0) circle (\de);
%   \end{scope}}
% }

\tikzset{ion/.pic={
  \begin{scope}[transparency group,nearly transparent]
  \begin{scope}[transparency group,semitransparent]
  \shade[ball color=red] ([shift={(rand*180:rnd*\de cm)}]0,0) circle (\dn);
  \end{scope}
  \begin{scope}[transparency group]
  \shade[ball color=red] ([shift={(rand*180:{rnd*\de cm})}]0,0) circle (\dn);
  \end{scope}
  \end{scope}
  \shade[ball color=red] ([shift={(rand*180:{rnd*\de cm})}]0,0) circle (\dn);}
}

\tikzset{es/.pic={
  \begin{scope}[transparency group,nearly transparent]
  \begin{scope}[transparency group,semitransparent]
  \shade[ball color=NavyBlue!70] ([shift={(180:2*\de cm)}]0,0) circle (\de);
  \end{scope}
  \begin{scope}[transparency group]
  \shade[ball color=NavyBlue!70] ([shift={(180:{\de cm})}]0,0) circle (\de);
  \end{scope}
  \end{scope}
  \shade[ball color=NavyBlue!70] (0,0) circle (\de);}
}

\def\dn{.15}
\def\de{.1}


%atoms left
\foreach \y in {0,...,1}{
\foreach \x in {0,...,4}{
\pic[] at ({\x},{\y}) {ion};
\pic[] at ([shift={(rnd*360:{.4 cm})}]\x,\y) {es};
}}


    \end{tikzpicture}
\caption{Направленное движение электронов между ионами металла}
\label{fig:p120}
\end{figure}

Проводник, обладающий сопротивлением, называют \emph{резистором} и изображают на схемах следующим образом (рис.\,\ref{fig:p121}).

    \begin{figure}[H]
    \centering
\begin{circuitikz}[font=\small,thin,
    line cap=round,line join=round,
>={Stealth[inset=0pt,round,length=3mm, angle'=20]},
vec/.style={>={Stealth[inset=0pt,round,length=3.5mm, angle'=20]}}]

    \ctikzset{bipoles/americaninductor/coils=3}

    \draw (0,1) to[*R=$R$] ++ (2,0);

\end{circuitikz}
\caption{Резистор}
\label{fig:p121}
\end{figure}



\end{NoHyper}
\end{document}
